\section{Introduction}\label{sec:intro120}
An economic decision rule must choose a feasible action using the information
available at the decision date. Its numerical value is determined by the
resulting recursive economic objective and by the resources needed to deliver
that action. Neural Bellman Operators (NBOs) study the construction and learning
of continuation operators for this purpose. A learned critic represents the
continuation entering a Bellman problem; an actor obtains an admissible decision
from that represented problem; and the value of the implemented policy follows
its own future decisions under the original economic recursion. This distinction
becomes consequential with nonlinear recursive risk, state-dependent action
constraints, imperfect optimization, and finite information acquisition.

The central question is when learning a continuation operator improves economic
decisions at a competitive complete cost. Answering it requires more than a
function-approximation result. A small fitted value residual can coexist with a
poor action direction, an incorrect boundary decision, an expensive risk
evaluation, or an uncertified future policy. A representation that is cheap to
query can also be expensive to acquire or to qualify. The original NBO framework
therefore keeps the target Bellman recursion, the implemented policy recursion,
and the computational account explicit. This revision advances that framework
through action comparisons whose information requirements depend on the actions
being compared, followed by a new original-objective policy-improvement study.

Our principal application is the original constrained investment economy.
A firm has a vector of productive states, a state-dependent investment capacity,
a nonlinear propagation technology, and continuous innovations. Current and
terminal costs penalize distance from productive targets, cross-component
imbalances, and severe aggregate shortfalls. Upper-tail cost risk is evaluated
recursively, so the adverse continuation distribution affects the investment
decision at each date. The model has a uniquely defined original Bellman optimum and a
measurable implemented terminal controller. These objects permit economic
comparisons of neural, hinge, polynomial, spline, grid and semi-analytic
continuations using one original objective and feasible action set, with common
access to the economic primitives. The application is a controlled economic
numerical experiment; its cost
units and initial-state distribution are declared rather than estimated from
empirical firm data.

\paragraph{An action-scaled Bellman comparison.}
The first contribution is a quantitative comparison modulus derived from the
original constrained technology and recursive preference. Let
$f_a(z)=V_1(F(x,a,z))$ denote the original optimal continuation after action $a$.
For the two-dimensional investment problem we prove a global Lipschitz bound
$M_\alpha$ for the state gradient of $V_1$, including binding investment
constraints and changes of the active risk tail. Normalizing the feasible action
interval makes the boundary calculation explicit. It follows that
$\operatorname{Lip}(f_b-f_a)\le(4/5)M_\alpha|b-a|$.
Thus the fixed-density integration error of the continuation difference admits
a bound proportional to action distance, even when the level of each
continuation remains substantial.

This modulus gives a computable paired midpoint interval. Intersecting it with
the separately valid convex midpoint--secant intervals yields the integration
width
\[
 \min\{2KR(1-\alpha)n^{-2},\ E_{ab}R(1-\alpha)n^{-1}\},
 \qquad E_{ab}=(4/5)M_\alpha|b-a|,
\]
under a known feasible risk-tail density and exact point information. Acquired
node errors, numerical accumulation, tail localization and future-policy error
are then charged explicitly. We prove a matching worst-case order for a precise
fixed-density scalar-value information interface. The comparison applies to all
representations given that interface. The additional smoothness and algebraic
access available in the investment economy are stated separately; the abstract
lower bound does not remove these advantages from a conventional solver.

The second contribution connects action alignment to the movement of the
adverse tail. A common feasible density has a nonnegative deficit relative to
the maximizing density of each continuation. Strong convexity in the innovation
and the action-scaled comparison modulus transfer a bound for one deficit to
the other, including a tail optimizer at the boundary. At an exactly optimal
candidate tail, the resulting mismatch allowance is quadratic in action
distance. The implemented comparison keeps both directly acquired deficits and
intersects the incumbent deficit with this transfer bound. It retains the
original risk supremum throughout.

\paragraph{A policy change must improve the implemented policy.}
The new gate starts from a completed, qualified incumbent. A candidate is
accepted only when its original complete-policy contrast has a strictly negative
upper bound and the candidate satisfies the unchanged original regret and
interval-width requirements. Every accepted step therefore has a certified
economic gain. Every refusal, unresolved comparison, and proposal acquisition
remains in the cost account. The three projected support proposals and the
completed-incumbent safeguard are retained methods; the new numerical mechanism
changes the information used to authorize their realized policy changes.

We prove a finite-resolution sufficient condition for acceptance, with distinct
allowances for common-tail mismatch, point accuracy, integration, and actual
future behavior. There is also a sharp limit to the stated information: knowing
the optimal-future action values exactly leaves a comparison uncertainty of
$\pm\beta\varepsilon_1$ when the implemented future policy is known only to
have error in $[0,\varepsilon_1]$. Resolving smaller gains requires information
about that actual future policy. This result identifies a specific numerical
research problem within NBO rather than treating an unresolved comparison as
optimality or interpreting additional internal fitting accuracy as economic
progress.

\paragraph{Evidence in the original economic units.}
The new experiment fixes eight continuation representations and four complete
services before independent initial states are generated. In addition to the
completed incumbent and the new paired service, it includes the retained
three-trial safe method and a stronger direct original-continuation proposal
that receives the same paid comparison gate. Every learned representation uses
the same original $V_1$ value labels. The hinge control receives the neural
features and pays their acquisition in full. The semi-analytic original solver,
quadratic and polynomial controls, coordinate spline, and conventional grid
remain eligible competitors. Actual information counts and complete costs are
reported even when they differ across methods.

The independently generated 64 state rows support 4,096 primary decisions.
Three additional serial processes repeat a fixed 32-state subset for descriptive
timing, without creating additional independent observations. A separate exact
rational implementation reconstructs every unique returned state/action/risk
key. It evaluates the original continuous terminal risk and also constructs an
all-action lower bound for the original date-zero optimum. This gives original
policy regret and action-distance bounds in addition to fixed-policy value
enclosures. Prespecified binomial and bounded-mean inference transfer certified
gains over each service's own incumbent to the declared state population.
Cross-representation adoption still requires comparing complete economic costs
and all acquisition and decision-service resources.

\paragraph{Relation to numerical control and learning.}
Monotone improvement from approximate policy evaluation has a substantial
literature. \citet{metelli2021} construct policy-improvement lower bounds and safe
policy-iteration schemes. \citet{eshwar2026} study Bellman-constrained evaluation
and conservative improvement with function approximation. These works make the
general need for an improvement guarantee clear. The present addition concerns
the paid continuous-risk action-comparison problem generated by the original
NBO economy, its explicit constrained smoothness and tail-mismatch constants,
and the policy consequences of that information.
Accuracy certificates for convex minimization with inexact oracles also supply
valid stopping criteria \citep{gladin2023}. Here the economic recursion requires
separately accounting for continuation, risk, feasible actions, and the actual
future controller; an internal optimization residual alone does not provide
those accounts. Classical binomial inversion and concentration are attributed
when used, and ordinary hinge or polynomial expressiveness is not credited to
neural networks alone.

The main article presents the NBO construction, original recursive-risk account,
investment model, new action comparison and economic evidence. The Technical
Supplement proves the Bellman foundations, stability, own-future errors,
constrained gradients, full critic--actor graph, information statements and
new results. It also retains the separate Brownian recursive-utility and stopped
boundary bridge under its own assumptions. The Development Appendix preserves
the broader economic program, previous theorems, complete historical comparisons,
failed or adverse results, and source-specific responses. Historical experiments
retain their own protocols and evidence identities; the new independent sample
is used only for its prospectively declared study.
